\documentclass[12pt, titlepage]{article}
\usepackage[margin=1in]{geometry}
\usepackage{graphicx}
\usepackage{float}

\usepackage{hyperref}

\title{\Huge Wingfeather Flight Controls Lab Drone \\ \Large{User Manual}}
\author{by \\ \Large TEAM WINGFEATHER \\ Aydan Vissing \\ Brady Schick \\ Gabriel Southwell \\ Simon Ross}
\date{\footnotesize{Last Updated} \\ \normalsize{\today}}


%\setlength{\topmargin}{-.25in}
%\setlength{\oddsidemargin}{.25in}
%\setlength{\evensidemargin}{.25in}
%\setlength{\textwidth}{6.0in}
%\setlength{\textheight}{9.0in}
%\renewcommand{\baselinestretch}{1.4}
%\usepackage{indentfirst} %Maybe keep this maybe not

\begin{document}
\maketitle

%FIXME remove comments

%FIXME improve abstract

\newpage
\tableofcontents
%\listoffigures
%\listoftables

\newpage
\section{Quick Start Guide}

This drone is designed to be used in-lab with just the Wingflasher application and the required hardware. Assembling the hardware can be challenging 

The user manual must include a Quick Start guide which should be a very simple one-page explanation of how to get the device set up for a first test using built-in default settings.

You must explain how to use all your instrument’s features. Assume that the reader of your user’s manual has never seen such an instrument. Your user manual will be their only reference source as they try to set up (i.e. connect the wires) and make use of the device. You may want to test your user manual on someone who has no experience like a friend from a different major. You may assume that the general user has a basic understanding of the terminology common to the project.

\section{Drone Construction}

\subsection{Ordering Components}
\subsubsection{PCB}

Essential files for working with the Wingfeather circuit board are located in the \href{https://github.com/Team-Wingfeather/fcl_senior_design/tree/main/pcb}{\texttt{pcb} directory}. This includes both the editable EasyEDA files and Gerber files for ordering. 

Ordering printed circuit boards is made simple by \href{https://jlcpcb.com/}{JLCPCB}, which is also directly accessible from an EasyEDA project. In general, the cheapest settings can be used for ordering.

\subsubsection{Motors}
The current motors used by the Wingfeather project are 720 coreless brushed DC motors. Most of these found on Amazon are about the right size, but stock tends to fluctuate. The current ones being used are found \href{https://www.amazon.com/dp/B0BBR16BK3}{here}.

\subsubsection{Propellers}
(key specs and a link)

\subsection{Printing the Frame}

\subsection{Assembling the components}

Once the necessary components have been obtained, assembling the drone is fairly straightforward. 
\begin{enumerate}
\item Solder the battery connector of choice to the bottom of the board. The plastic part of the connector can slide off of the metal part, so it can be helpful to super-glue the plastic to the PCB before soldering. Note that there are two connector options: JST-PH and JST-XH. To date, the team has been using JST-PH connectors for their minimal profile. Be careful to match the polarity of the batteries to the polarity of the board, marked with $+/-$ symbols near the connector.
\item Using header pins, solder 2 time-of-flight sensor daughter boards onto the main board---one on top and the other on the bottom. These should use the 6-pin header slots, and be soldered onto the face with the printed labels.
\item Place the PCB on the frame and use 7 M2 screws to connect them. Note that the orientation of the PCB is determined by which corner does not have a hole for a screw. The front of the drone is the side with the USB connection.
\item Connect the propellers onto the motors. Make sure to select where each motor will go and match the propeller direction with the direction indicated around the screw nearest to that motor. Pushing gently onto the table can help affix the propellers to the shafts. Finally, it  is conventional to use a different color of propellers to designate the front of the drone, depending on the model purchased.
\item Slide the 4 motors into the designated recesses. It may prove easier to solder the motor connections before gluing the motors in.
\item Solder each motor to the nearest 2 through holes on the circuit board. Note that polarities of the annular rings alternate around the drone. This is so that the motors across from each other spin the same direction and adjacent motors spin opposite directions. In practical terms, the square annular ring should always be connected to the same motor terminal ``polarity'' (designated by the color of the wire). Which color may need to be determined by trial and error. For the current motors, the white motor lead should be soldered to the square annular ring. In the end, make sure the direction each motor spins matches the arrow located around the nearest PCB screw. Note that it may be helpful to twist the wires from each motor before soldering. After soldering, this twisted pair may be hot glued to the frame to save space. Use caution as the soldering iron can easily melt the 3D-printed frame.
\item Slide each motor partially out and apply a moderate amount of hot glue under it. Slide the motor back down, and be ready to wipe away excess glue. Affixing the motors well is important so they do not fly off in flight. Be careful, as the glue gun can melt the plastic of the frame.
\item In the future, more time-of-flight sensors may be added to the 5-pin header slots. Additionally, a velcro strap can be used to affix the battery to the bottom of the board through the designated slot.
\end{enumerate}


\subsection{Programming the Drone}

\subsection{Using the App}

\subsection{Using ESP IDF}

\section{Drone Development}

\subsection{Basic Repository Structure}
Don't look at the commit history.

\subsection{Basic Drone Testing}

\subsection{PID Tuning}

\subsection{Frame Development}

\subsection{PCB Development}
Printed circuit board development can be accomplished easily using \href{https://easyeda.com/}{EasyEDA}. This software allows for design of a schematic and then layout of the circuit board. There is also a web interface that works quite well.

The current board revision files can be found \href{https://github.com/Team-Wingfeather/fcl_senior_design/tree/main/pcb/v1/EasyEDA}{here}. The zipped project folder contains all of the past board revisions, but the most up-to-date is  \texttt{PCB\_FCL\_Drone\_v0.0.13\_FAB}.
The current board revision is out of date in a few areas: 
\begin{itemize}
\item The current power regulator (\href{https://github.com/Team-Wingfeather/fcl_senior_design/blob/main/resources/datasheets/C109322-TPS63070RNMR.pdf}{TPS63070}) is weak and burns out. This component will need to be replaced in the next version of the board, potentially by a \href{https://jlcpcb.com/partdetail/HoltekSemicon-HT75331/C14289}{HT7533-1} LDO. This is comparable to what is currently soldered onto the board. Note: there are multiple capacitors, resistors, and inductors that are used only by the previous component and will also need to be removed. See the \href{https://github.com/Team-Wingfeather/fcl_senior_design/blob/main/resources/datasheets/C109322-TPS63070RNMR.pdf}{datasheet} and \href{https://github.com/Team-Wingfeather/fcl_senior_design/tree/main/pcb/v1/EasyEDA}{EasyEDA} files for which components to replace.
\item The bottom ToF sensor leaves little room for a velcro strap, so it could be moved, or the frame CAD file could be modified to accommodate.
\item Resistor 13 on the board has been halved by a soldered-on parallel resistor of equal value. This creates a 1/3 voltage divider, bringing the 2C 7.4V battery voltage into a readable range for the ESP32. Thus, on board version 2, R13 should be replaced with a 50K Ohm resistor.
\end{itemize}
Some other optional changes include:
\begin{itemize}
\item The green LED is painfully bright and could use a larger resistor to dim it.
\item Connecting the ToF interrupt pins to the ESP32 for marginally faster sensor response.
\item Moving the MPU6050 IMU to its own I2C line.
\item Switching GPIO 23 to GPIO 16. The motor 4 PWM pin failed in testing and was swapped to pin 16. This should likely be made a permanent hardware change. See \href{https://github.com/Team-Wingfeather/fcl_senior_design/commit/29bd241f75058da32211addb992e50635beb5c23}{this commit} that patches the software.
\end{itemize}

\subsection{App Development}



\end{document}
